\chapter{软开关技术}

现代电力电子装置的发展方向是高频化、小型化、高功率密度。但开关频率越高，\textbf{开关损耗}越大，器件发热、效率下降、电磁干扰加剧——这个矛盾靠“硬开关”本身无法解决。\textbf{软开关技术}就是在开关过程前后引入谐振（或辅助电路），使开关在\textbf{电压为零}或\textbf{电流为零}的时刻动作，从而把开关损耗降到接近零。

% ============================================================
\section{软开关的基本概念}

\subsection{硬开关与软开关}

前面各章的分析都把开关理想化，认为器件瞬间完成通断。实际上器件的开通和关断都需要时间，在这段时间内\textbf{电压和电流同时不为零}，其乘积即为损耗。

\begin{itemize}
  \item \textbf{开通损耗}：开通时电流上升、电压下降，二者交叠产生损耗；若器件两端并联电容，开通瞬间电容储能经器件释放，还会产生很大的\textbf{电流尖峰}（容性开通）。
  \item \textbf{关断损耗}：关断时电流下降、电压上升，二者交叠产生损耗；若回路中有电感，关断瞬间电感储能产生很高的\textbf{电压尖峰}（感性关断），可能击穿器件。
  \item 此外，硬开关还伴随严峻的\textbf{电磁干扰}（du/dt、di/dt 很大）与噪声。
\end{itemize}

\begin{tcolorbox}[enhanced,breakable,colback=highlight,colframe=yhbAccent!30!white,boxrule=0.5pt,arc=2.5pt,left=11pt,right=11pt,top=9pt,bottom=9pt,boxsep=0pt,before skip=11pt,after skip=11pt]
\textbf{硬开关}\quad 开关过程中电压、电流均不为零，出现了\textbf{交叠}，因而有开关损耗。\\
\textbf{软开关}\quad 通过谐振等手段，使开关\textbf{开通前电压先降到零}，或\textbf{关断前电流先降到零}，从而消除电压电流的交叠，使开关损耗为零（理论上）。软开关同时还能抑制 du/dt、di/dt，降低电磁干扰。
\end{tcolorbox}

\subsection{零电压开关与零电流开关}

\begin{itemize}
  \item \textbf{零电压开关（ZVS）}：使开关开通前其两端电压为零，则开通时不会产生损耗（也就没有容性开通电流尖峰）。要实现 ZVS，须让开关两端的电压在开通前先谐振到零，并让反并联二极管先导通（把电压钳在零）再给开关开通信号。
  \item \textbf{零电流开关（ZCS）}：使开关关断前其电流为零，则关断时不会产生损耗（也就没有感性关断电压尖峰）。要实现 ZCS，须让流过开关的电流在关断前先谐振到零，并让反并联二极管先导通。
\end{itemize}

\corepoint{记住对应关系：\textbf{ZVS 对应“开通”（电压为零）}、\textbf{ZCS 对应“关断”（电流为零）}；实际电路中往往同时安排“零电压开通 + 零电流关断”（如移相全桥），即所谓“全软开关”。判断某电路属于哪种软开关，就看\textbf{谐振元件（$L_r$、$C_r$）与开关的连接方式}以及\textbf{开关动作时刻电压电流是否为零}。}

% ============================================================
\section{软开关电路的分类}

按实现软开关所用电抗元件与开关的连接方式，软开关电路分为三大类。

\subsection{准谐振电路}

在开关两端（或串联）接入 $L_r$、$C_r$ 构成谐振回路，使电压或电流在开关动作时刻形成\textbf{正弦半波}（故称“准谐振”）。分为：

\begin{itemize}
  \item \textbf{零电压开关准谐振电路（ZVS-QR）}：$L_r$、$C_r$ 串联后并联在开关两端，开关开通前 $C_r$ 上的电压先谐振到零；
  \item \textbf{零电流开关准谐振电路（ZCS-QR）}：$L_r$、$C_r$ 与开关串联，开关关断前 $L_r$ 中的电流先谐振到零；
  \item \textbf{多谐振电路（MRC）}：电路中存在多个谐振元件，可同时实现开关与二极管的软开关。
\end{itemize}

谐振频率与特性阻抗决定了谐振过程的快慢与峰值：

\keyformula{\omega_r=\frac{1}{\sqrt{L_rC_r}},\qquad Z_r=\sqrt{\frac{L_r}{C_r}}}

准谐振电路的共同缺点是：\textbf{开关应力大}（谐振电压/电流峰值高）、\textbf{输出与开关频率相关}（只能靠调频调压，不利于滤波设计），且软开关条件与负载有关（轻载时可能失去软开关）。

\subsection{零开关 PWM 电路}

在准谐振电路基础上加辅助开关，使电路\textbf{只在开关动作前后的短时间内谐振}，其余时间按 PWM 方式工作。分为\textbf{零电压开关 PWM 电路（ZVS-PWM）}与\textbf{零电流开关 PWM 电路（ZCS-PWM）}。

优点：既实现了软开关，又保持了 PWM 控制（恒频、调宽）的优点，输出电压与占空比成线性关系。

\subsection{零转换 PWM 电路}

谐振网络\textbf{与主开关并联}（而非串联），且只在开关动作时投入工作，主电路电流不流过谐振网络，因此\textbf{导通损耗小、软开关范围宽（轻载也能实现）}。分为\textbf{零电压转换 PWM 电路（ZVT-PWM）}与\textbf{零电流转换 PWM 电路（ZCT-PWM）}。

\begin{table}[htbp]
\centering
\caption{三类软开关电路的比较}
\label{tab:8-1}
\begin{tabular}{p{3.2cm}p{4.2cm}p{6.6cm}}
\toprule
\textbf{类型} & \textbf{谐振元件位置} & \textbf{特点} \\
\midrule
准谐振（QR） & 与开关串联或并联 & 谐振贯穿整个开关周期，应力大、调频控制、软开关范围受负载影响 \\
零开关 PWM（ZS-PWM） & 与开关串联或并联，加辅助开关 & 只在开关瞬间谐振，恒频 PWM，应力仍较大 \\
零转换 PWM（ZT-PWM） & 与主开关并联，辅助支路 & 主电路电流不经过谐振网络，损耗小、软开关范围宽，是主流方案 \\
\bottomrule
\end{tabular}
\end{table}

% ============================================================
\section{典型的软开关电路}

\subsection{零电压开关准谐振电路}

这是结构较为简单的软开关电路，以 Buck 型为例：在开关 $S$ 两端并联谐振电容 $C_r$，在主回路串入谐振电感 $L_r$。

\begin{itemize}
  \item $S$ 关断时，$L_r$ 中的电流向 $C_r$ 充电，$S$ 两端电压 $u_S$ 从零开始上升（\textbf{零电压关断}，无容性关断损耗）；随后 $L_r$、$C_r$ 谐振，$u_S$ 形成正弦半波；
  \item 当 $u_S$ 谐振回到零时，反并联二极管先导通把 $u_S$ 钳在零，此时再给 $S$ 开通信号，即实现\textbf{零电压开通}；
  \item 谐振电压峰值约为 $2E$，因此\textbf{开关管耐压须按 $2E$ 以上选取}，这是 ZVS-QR 的主要缺点；此外输出电压与开关频率成反比（调频调压），轻载时可能失去 ZVS。
\end{itemize}

\subsection{谐振直流环}

谐振直流环（Resonant DC Link，RDCL）适用于变频器：在直流母线与逆变桥之间串入小电感 $L_r$，在逆变桥的直流侧并联电容 $C_r$，使\textbf{直流母线电压周期性谐振到零}，从而为逆变桥的所有开关提供零电压开通条件。

\begin{itemize}
  \item 母线电压每谐振到零一次，逆变桥就可以“择机”进行一次开关动作（离散脉冲调制），从而实现全桥 ZVS；
  \item 优点：所有开关都能零电压开通，控制灵活；
  \item 缺点：母线电压峰值高（约 $2E$ 以上）、器件电压应力大，且谐振过程增加了额外的导通损耗与 EMI。
\end{itemize}

\subsection{移相全桥型零电压开关 PWM 电路}

这是目前中大功率 DC-DC 变换器最常用的软开关拓扑：全桥的四个开关中，\textbf{超前桥臂}与\textbf{滞后桥臂}的驱动信号之间加入\textbf{移相角}，利用变压器\textbf{漏感}（或外加谐振电感）与开关管的\textbf{结电容}谐振实现 ZVS。

\begin{itemize}
  \item 每个桥臂的上下管互补导通（各约 50\%），两桥臂之间移相，移相角决定输出电压（占空比）；
  \item \textbf{死区时间内}，原导通的管子关断后，负载电流经另一管的结电容与漏感谐振，把即将开通的管子两端电压谐振到零，再由其反并联二极管钳位，从而实现 ZVS；
  \item 特点：\textbf{超前桥臂容易实现 ZVS}（利用滤波电感能量），\textbf{滞后桥臂较难}（只靠漏感能量，轻载时可能失去 ZVS）；
  \item 存在\textbf{占空比丢失}：换相期间二次侧电压被“短路”，有效占空比小于一次侧给定占空比，输出电压比理论值低。
\end{itemize}

\subsection{零电压转换 PWM 电路}

零电压转换（ZVT-PWM）电路在主开关两端并联一个由辅助开关 $S_a$、谐振电感 $L_r$、谐振电容 $C_r$ 组成的辅助支路：

\begin{itemize}
  \item 主开关开通前，先开通辅助开关 $S_a$，$L_r$ 中的电流线性上升，把主开关的电流“抽走”，使主开关反并联二极管导通、电压钳在零；
  \item 然后开通主开关，即为\textbf{零电压开通}；随后关断 $S_a$，$L_r$ 的能量经二极管回馈，辅助开关也可实现软关断；
  \item 优点：\textbf{辅助支路只在开关瞬间工作}，主电路电流不经过 $L_r$，导通损耗小；\textbf{软开关范围宽}，轻载仍能实现 ZVS；辅助器件容量小、成本低。
\end{itemize}

% ============================================================
\section{软开关技术新进展}

软开关技术的发展方向主要有：

\begin{enumerate}
  \item \textbf{无源无损缓冲电路}：只用电感、电容、二极管构成缓冲网络，不加辅助开关，靠谐振把开关损耗转移到缓冲元件中再回馈，结构简单、可靠性高，但软开关范围有限。
  \item \textbf{有源钳位电路}：用一只辅助开关加钳位电容，把开关两端的电压尖峰钳位并回收能量，同时实现主开关的 ZVS，广泛用于正激、反激等单端电路。
  \item \textbf{LLC 谐振变换器}：用谐振网络（$L_r$、$C_r$ 与励磁电感 $L_m$）实现全负载范围的 ZVS 与二次侧 ZCS，效率高、EMI 小，是目前服务器电源、充电器的热门拓扑。
  \item \textbf{复合软开关与多电平软开关}：把软开关技术与多电平、多重化技术结合，用于高压大容量场合；也采用“几个简单高效的开关电路级联、并联、串并联组合”的思路，用成熟单元构成高性能系统。
  \item \textbf{软开关与数字控制结合}：用数字信号处理器实时调整死区与移相角，保证全工况下的软开关条件（自适应死区控制）。
\end{enumerate}

\corepoint{本章的考点是“概念辨析 + 分类归属”：\textbf{（1） 硬开关损耗来自电压电流交叠，软开关靠谐振消除交叠}；\textbf{（2） ZVS 对应开通、ZCS 对应关断}；\textbf{（3） 三类电路的归属}——准谐振（谐振贯穿全周期、应力大、调频）→ 零开关 PWM（只瞬间谐振、恒频、应力仍大）→ 零转换 PWM（谐振网络与主开关并联、损耗小、范围宽）；\textbf{（4） 典型电路的软开关机理}：ZVS-QR 靠 $L_rC_r$ 谐振到零、谐振直流环靠母线电压周期性到零、移相全桥靠漏感与结电容谐振（滞后桥臂轻载难 ZVS、有占空比丢失）、ZVT 靠辅助支路先抽走电流。}
